\documentclass[12pt]{article}
\usepackage[slovene]{babel}

\usepackage{ragged2e,lmodern} %sodobni LaTeX

\usepackage{array,booktabs, multirow} % Tabele
\usepackage{graphicx} % Slike

\usepackage{mathtools}
\usepackage{amsthm, amssymb} % Matematika
\usepackage{lipsum}

\usepackage{hyperref} % Hiperpovezave



\title{\LaTeX{} 4. Matematika (izreki, definicije, zgledi)}
\author{Antonio Montero}
\date{7. november 2025}



\begin{document}
\maketitle



  % \begin{izrek}[Pitagorov izrek]
  %   Vsota površin kvadratov katet pravokotnega 
  %   trikotnika je enaka površini kvadrata nad 
  %   hipotenuzo.
  %   Izrek lahko zapišemo tudi kot: \[c^{2} = a^{2} + b^{2}\]
  % \end{izrek}

  % \begin{izrek}[Murphyjev zakon]
  %   Če je nekaj lahko narobe, bo šlo narobe.
  % \end{izrek}

  % \begin{izrek}[Homer Simpson]
  %   Kljub Fermatovi zadnji teoremi velja naslednja enačba: \[1782^{12} + 1841^{12} = 1922^{12}\]
  % \end{izrek} 

  
\end{document}